\clearpage
\fancyhead[L]{\small COMET.}
\SourceLeaf{422}
\begin{center}
\Large COMET.

\medskip
Göttingische gelehrte Anzeigen. Stück 28. Seite 273..278. 1819 Febr. 18.
\end{center}

\par\indent Ueber die beiden von \textsc{Pons} in Marseille im November des vorigen Jahrs entdeckten
Cometen theilen wir hier einige Nachrichten um so lieber mit, da diese Cometen in
mehreren Beziehungen zu den vorzüglich merkwürdigen gehören. Der eine scheint in
Deutschland zuerst auf der Mannheimer Sternwarte beobachtet zu sein. Die
Beobachtungen, welche daselbst von Hrn. Prof. \textsc{Nicolai} angestellt und in einem
Schreiben an den Hrn. Hofrath \textsc{Gauss} vom 8. Januar mitgetheilt sind, waren folgende:

\begin{center}
\setlength{\tabcolsep}{2.2pt}\begin{tabular}{c|ccc|ccc|ccc}
1818 & \multicolumn{3}{c|}{\shortstack{Mittlere Zeit\\in Mannheim}} & \multicolumn{3}{c|}{\shortstack{Gerade\\Aufsteigung}} & \multicolumn{3}{c}{\shortstack{Nordliche\\Abweichung}} \\
\hline
Dec. 22 & $6^{\mathrm u}$ & $51^{\mathrm m}$ & $48^{\mathrm s}$ & $326^\circ$ & $18'$ & $13''$ & $2^\circ$ & $54'$ & $5''$ \\
        & 8     & 45     & 37     & 326         & 16    & 59     & 2         & 52    & 24    \\
     23 & 7     & 11     & 25     & 326         & 3     & 39     & 2         & 40    & 26    \\
     24 & 7     & 14     & 3      & 325         & 48    & 33     & 2         & 26    & 59    \\
     25 & 6     & 48     & 49     & 325         & 32    & 49     & 2         & 13    & 34    \\
     29 & 7     & 0      & 31     & 324         & 19    & 58     & 1         & 10    & 57    \\
\end{tabular}
\end{center}

\Needspace{10\baselineskip}\par\indent Zugleich waren folgende parabolische Elemente beigefügt, welche Hr. Prof. \textsc{Nicolai}
vorläufig aus den Beobachtungen vom 22., 25. und 29. December berechnet hatte.

\begin{center}
\begin{tabular}{l l}
Zeit der Sonnennähe 1819 Jan. . . . . . . . & 25.0000 in M. \\
Logarithm des kleinsten Abstandes . . . . . . & 0.52933 \\
Länge der Sonnennähe . . . . . . . . . . & $145^\circ$ $49'$ 10 \\
Länge des aufsteigenden Knoten . . . . . . . & 330 \quad 14 \quad 17 \\
Neigung der Bahn . . . . . . . . . . . . & 14 \quad 59 \quad 6 \\
\multicolumn{2}{l}{Richtung: \emph{rechtläufig.}} \\
\end{tabular}
\end{center}

\Needspace{14\baselineskip}\par\indent Auf der hiesigen Sternwarte wurden durch Hrn. Prof. \textsc{Harding} folgende Beobachtungen
gemacht:

\begin{center}
\setlength{\tabcolsep}{2.2pt}\begin{tabular}{c|ccc|ccc|ccc l}
& \multicolumn{3}{c|}{\shortstack{Mittlere Zeit\\in Göttingen}} & \multicolumn{3}{c|}{\shortstack{Gerade\\Aufsteigung}} & \multicolumn{3}{c}{Abweichung} & \\
\hline
1818 Dec. 25 & $6^{\mathrm u}$ & $39^{\mathrm m}$ & $52^{\mathrm s}$ & $325^\circ$ & $33'$ & $30''$ & $2^\circ$ & $14'$ & $49''$ & Nordl. \\
          26 & 5     & 58     & 32     & 325         & 17    & 8      & 2         & 0     & 31     & \\
          27 & 6     & 7      & 10     & 324         & 59    & 15     & 1         & 45    & 5      & \\
          28 & 6     & 24     & 8      & 324         & 40    & 17     & 1         & 27    & 31     & \\
1819 Jan.  1 & 6     & 36     & 16     & 323         & 11    & 48     & 0         & 14    & 58     & \\
           7 & 6     & 39     & 41     & 319         & 53    & 59     & 2         & 19    & 59     & Südl. \\
           8 & 7     & 37     & 34     & 319         & 15    & 28     & 2         & 51    & 49     & \\
          12 & 6     & 6      & 7      & 315         & 36    & 53     & 5         & 35    & 24     & \\
\end{tabular}
\end{center}

\Needspace{11\baselineskip}\par\indent Auf der Seeberger Sternwarte beobachtete Hr. Prof. \textsc{Encke} diesen Cometen fünfmal:

\begin{center}
\setlength{\tabcolsep}{2.2pt}\begin{tabular}{c|ccc|ccc|ccc l}
& \multicolumn{3}{c|}{\shortstack{Mittlere Zeit\\in Seeberg}} & \multicolumn{3}{c|}{\shortstack{Gerade\\Aufsteigung}} & \multicolumn{3}{c}{Abweichung} & \\
\hline
1819 Jan. 1 & $6^{\mathrm u}$ & $38^{\mathrm m}$ & $46^{\mathrm s}$ & $323^\circ$ & $11'$ & $36''$ & $0^\circ$ & $14'$ & $38''$ & Nördl. \\
          4 & 7     & 26     & 56     & 321         & 43    & 43     & 0         & 54    & 59     & Südl. \\
          5 & 7     & 12     & 12     & 321         & 10    & 3      & 1         & 20    & 57     & \\
          6 & 5     & 49     & 48     & 320         & 35    & 14     & 1         & 47    & 57     & \\
         12 & 6     & 16     & 32     & 315         & 35    & 20     & 5         & 35    & 26     & \\
\end{tabular}
\end{center}

\SourceLeaf{423}
\expandafter\def\csname HistoricalNoteText423\endcsname{wahrscheinlich Seeberger Zeit, was von Hrn. \textsc{Encke} nicht ausdrücklich bemerkt war.}
\par\indent Die Resultate welche Hr. Prof. \textsc{Encke} aus seinen eignen und den Mannheimer
Beobachtungen herausgebracht, und Hrn. Hofr. \textsc{Gauss} in einem Schreiben vom 5. Febr.
mitgetheilt hat, verdienen die grösste Aufmerksamkeit. Obgleich diese Beobachtungen nur
einen Zeitraum von drei Wochen umfassen, so liessen sie sich doch nicht in einer Parabel
ohne Fehler von mehreren Minuten darstellen, Fehler, die bei diesen Beobachtungen nicht
zulässig waren. Nach mehreren Versuchen fand Hr. \textsc{Encke}, unter Zuziehung der freilich nur
sehr rohen \textsc{Pons}schen Angaben, eine elliptische Bahn von 3 Jahren und 7 Monaten
Umlaufszeit, welche alle Beobachtungen auf das trefflichste vereinigte.

\Needspace{16\baselineskip}\par\indent Folgendes ist die Uebersicht der Unterschiede der in Mannheim und Seeberg beobachteten
Oerter von den berechneten:

\begin{center}
\begin{tabular}{c|c|c|l}
& Ger. Aufst. & Abweich. & Beobachter \\
\hline
Dec. 22 & $+\ 7''6$  & $-\ 4''7$  & Nicolai \\
     23 & $+\ 0.3$   & $+17.3$    & --- \\
     24 & $-\ 1.0$   & $+11.7$    & --- \\
     25 & $+11.2$    & $-\ 5.5$   & --- \\
     29 & $+14.6$    & $-\ 1.8$   & --- \\
Jan.  1 & $+15.6$    & $+14.8$    & Encke \\
      4 & $+26.5$    & $+27.2$    & --- \\
      5 & $+30.6$    & $+12.9$    & --- \\
      6 & $+36.1$    & $+23.8$    & --- \\
     12 & $+\ 7.1$   & $+\ 3.0$   & --- \\
\end{tabular}
\end{center}

\par\indent Ohne auf die \textsc{Pons}schen Angaben Rücksicht zu nehmen, würde sich eine noch merklich
bessere Uebereinstimmung haben erreichen lassen.

\Needspace{14\baselineskip}\par\indent Diese elliptischen Elemente sind folgende:

\begin{center}
\begin{tabular}{ll}
Zeit der Sonnennähe 1819. Januar . . & $27.13417$\footnotemark \\
Länge der Sonnennähe . . . . . & $156^\circ\;14'\;8''$ \\
Länge des Knotens . . . . . . . & $334\;18\;8$ \\
\multicolumn{2}{l}{(beide auf das mittlere Aequinoctium von 1819 bezogen)} \\
\end{tabular}\par\smallskip
\begin{tabular}{ll}
Neigung der Bahn . . . . . . . & $13^\circ\;42'\;30''$ \\
Logarithm der halben grossen Axe . & $0.3697758$ \\
Excentricitätswinkel . . . . . . & $58^\circ\;57'\;24''$ \\
\end{tabular}
\end{center}
\footnotetext{\csname HistoricalNoteText423\endcsname}

\par\indent Hoffentlich werden in Zukunft noch frühere \emph{gute} Beobachtungen aus Frankreich
oder Italien bekannt werden, wodurch die Ungewissheit der Bahnbestimmung in engere
Grenzen gebracht werden wird. Noch wichtiger würde es sein, wenn noch einige gute
Beobachtungen nach dem Durchgange durch die Sonnennähe in südlichen Gegenden gemacht
werden könnten. In unsern Breiten, wo der Comet sich nur noch wenig über den südlichen
Horizont erheben wird, ist dazu keine Hoffnung.

\Needspace{15\baselineskip}\par\indent Höchst merkwürdig ist die grosse Uebereinstimmung dieser Elemente mit denen des ersten
Cometen vom Jahr 1805. Zur Vergleichung setzen wir dieselben, so wie sie damals von Hrn.
Prof. \textsc{Bessel} in einer Parabel berechnet wurden, hier her, und bemerken nur noch, dass
schon damals die Vergleichung der Beobachtungen mit der Rechnung fast unverkennbare
Spuren einer Abweichung der Bahn jenes Cometen von der Parabel zeigte, die jedoch von
keinem Astronomen weiter verfolgt zu sein scheinen.

\begin{center}
\begin{tabular}{l l}
\multicolumn{2}{l}{Durchgang durch die Sonnennähe 1818 November} \\
\multicolumn{2}{l}{18.13782 Pariser Zeit} \\
Länge der Sonnennähe . . . . . . . . . . . . & $147^\circ$ $51'$ $28''$ \\
Länge des aufsteigenden Cometen . . . . . . . . & 344\ \ 37\ \ 19 \\
Neigung der Bahn . . . . . . . . . . . . . . & 15\ \ 36\ \ 36 \\
Logarithm des kleinsten Abstandes . . . . . . . & 9.57820 \\
\end{tabular}
\end{center}

\SourceLeaf{424}\noindent (Bei den \textsc{Encke}schen Elementen des Cometen von 1818. 1819 wird dieser Logarithm
9.52579). Gewiss verdient die Frage, ob der jetzige Comet seit dem Jahr 1805 solche
Störungen erlitten haben kann, dass die Unterschiede in den Elementen sich dadurch
erklären lassen, und ob demnach beide Cometen für identisch zu halten sind, eine
sorgfältige Untersuchung.

\par\indent Den andern von \textsc{Pons} entdeckten Cometen fand seinerseits auch Hr. Prof. \textsc{Bessel} am 22.
December auf der Königsberger Sternwarte auf.

\Needspace{17\baselineskip}\par\indent Folgende Beobachtungen theilte dieser Astronom in einem Schreiben vom 14. Januar dem
Hrn. Hofr. \textsc{Gauss} mit:

\begin{center}
\setlength{\tabcolsep}{2.2pt}\begin{tabular}{c|ccc|ccc|ccc}
& \multicolumn{3}{c|}{\shortstack{Mittlere Zeit\\in Königsberg}} & \multicolumn{3}{c|}{\shortstack{Gerade\\Aufsteigung}} & \multicolumn{3}{c}{\shortstack{Nordliche\\Abweichung}} \\
\hline
1818 Dec. 22 & $7^{\mathrm u}$ & $6^{\mathrm m}$  & $51^{\mathrm s}$ & $303^\circ$ & $1'$  & $21''7$ & $36^\circ$ & $48'$ & $20''2$ \\
             & 7     & 53     & 35     & 303         & 10    & 14.7    & 36         & 48    & 29.8    \\
             & 10    & 21     & 47     & 303         & 37    & 29.1    & 36         & 51    & $0.3\!:$ \\
          24 & 18    & 22     & 50     & 311         & 56    & $29.1\!:$ & \multicolumn{3}{c}{.\quad .\quad .\quad .} \\
          25 & 6     & 22     & 20     & 313         & 17    & 17.2    & 37         & 7     & 53.1    \\
          26 & 6     & 4      & 56     & 315         & 38    & 48.2    & 37         & 4     & 1.2     \\
             & 10    & 33     & 0      & 316         & 3     & 18.1    & 37         & 2     & 35.2    \\
          27 & 6     & 55     & 0      & 317         & 39    & 48.8    & 36         & 58    & 45.6    \\
          28 & 6     & 0      & 22     & 319         & 24    & $50.7\!:$ & 36       & 53    & $31.6\!:$ \\
1819 Jan.  1 & 11    & 11     & 4      & 324         & 39    & 12.3    & 36         & 22    & $34.2\!::$ \\
           2 & 6     & 5      & 6      & 325         & 22    & 3.4     & 36         & 15    & 54.3    \\
\end{tabular}
\end{center}

\Needspace{12\baselineskip}\par\indent Folgende von Hrn. Prof. \textsc{Bessel} auf die Beobachtungen vom 22. und 27. December und vom
2. Januar gegründete vorläufige parabolische Elemente waren diesem Schreiben
beigefügt:

\begin{center}
\begin{tabular}{l l}
\multicolumn{2}{l}{Zeit der Sonnenähe 1818 December 4.0968 Pariser Zeit.} \\
Aufsteigender Knoten . . . . . . . . . . . . & $90^\circ$ $7'$ $29''$ \\
Neigung der Bahn . . . . . . . . . . . . . & 117\ \ 19\ \ 10 \\
Abstand des Perihel vom Knoten . . . . . . . . & 347\ \ \ 0\ \ 24 \\
Logarithm des kleinsten Abstandes . . . . . . . & 9.928324 \\
\end{tabular}
\end{center}

\Needspace{14\baselineskip}\par\indent Nach Anleitung dieser Elemente wurde dieser Comet vom Hrn. Hofr. \textsc{Gauss} sogleich als
ein sehr blasser Nebelfleck aufgefunden; die Elemente gaben die gerade Aufsteigung
gegen 40, die Abweichung aber $10'$ zu gross. Hierauf wurden durch Hrn. Prof. \textsc{Harding}
noch folgende Beobachtungen auf der hiesigen Sternwarte angestellt:

\begin{center}
\setlength{\tabcolsep}{2.2pt}\begin{tabular}{c|ccc|ccc|ccc}
& \multicolumn{3}{c|}{\shortstack{Mittlere Zeit\\in Göttingen}} & \multicolumn{3}{c|}{\shortstack{Gerade\\Aufsteigung}} & \multicolumn{3}{c}{\shortstack{Nordliche\\Abweichung}} \\
\hline
1819 Jan. 26 & $7^{\mathrm u}$ & $19^{\mathrm m}$ & $5^{\mathrm s}$  & $335^\circ$ & $23'$ & $7''$  & $35^\circ$ & $18'$ & $29''$ \\
          27 & 7     & 2      & 35     & 335         & 36    & 4      & 35         & 18    & 50     \\
          28 & 7     & 22     & 16     & 335         & 50    & 38     & 35         & 20    & 58     \\
          29 & 7     & 27     & 56     & 336         & 4     & 38     & 35         & 22    & 32     \\
          30 & 6     & 54     & 10     & 336         & 16    & 33     & 35         & 23    & 31     \\
\end{tabular}
\end{center}

\par\indent Bei der immer zunehmenden Lichtschwäche dieses durch seine starke geocentrische
Bewegung seit der ersten Entdeckung merkwürdigen Cometen werden schwerlich noch
spätere Beobachtungen zu erwarten sein. Ob auch bei diesem Cometen die vorhandenen
Beobachtungen zur Bestimmung der Abweichung der Bahn von der Parabel hinreichend, und
welches die Resultate sein werden, wird eine künftige Untersuchung lehren.
